\section{Three Empirical Anchors}

The empirical role of the indexed sources is narrow and disciplined:
\begin{enumerate}[nosep]
  \item \textbf{Signatures of conscious state.} The 2024 Communications
        Biology source is used for one strong claim only: conscious conditions
        are associated with robust and generalizable dynamical
        structure-function signatures rather than being describable only in
        loose psychological language. It secures permission to talk about
        organized state signatures; it does not authorize direct measurement of
        free will or a one-to-one map from every behaviour to a neural label.
  \item \textbf{Connectivity and stability.} The 2025 Communications Biology
        source is used for a second strong claim: global conscious states are
        related to structured connectivity, local dynamics, and stability, so
        persistence is not mere noise but an architectural problem. It secures
        permission to treat persistence as structured; it does not prove that
        every local behavioural frame is a whole-brain state.
  \item \textbf{Transition dynamics.} The 2025 Nature Neuroscience sleep
        source is used for a third strong claim: entry into a new regime can
        be abrupt, patterned, and tractable in dynamical language rather than
        remaining purely descriptive. It secures permission to speak of
        parameter-sensitive transitions; it does not prove that all behavioural
        redirections share the same mechanism as sleep onset.
\end{enumerate}

\begin{discussion}
These sources do not prove the full framework-inertia model, and the chapter
should not pretend otherwise. Their value is more disciplined. They justify
treating state, stability, and transition as serious biological topics. Once
that floor is established, the rest of the manuscript can talk about entry
cost, regime persistence, and redirection without sounding like a private
productivity metaphor detached from the science of organized state changes.
\end{discussion}

\begin{remark}[Why these three anchors travel together]
Taken together, the three sources secure a minimal but useful chain: one anchor
establishes signatures, one establishes stability architecture, and one
establishes transition dynamics. The chapter needs all three because a state
language without persistence or transition would still be too thin to support
the manuscript's later control logic.
\end{remark}
